% !TeX spellcheck = de_DE

% ---

\documentclass[%
    % alignment is either standard, bodytohead or headtobody; default is standard
    alignment       = standard,%
    % whether to show some auxiliary rules; default is false
    visualize       = false,%
    % mainaccentcolor must be a colorname from (my-)xcolor; default is guardsmanred
    mainaccentcolor = guardsmanred,%
    % otheraccentcolor must be a colorname from (my-)xcolor; default is guardsmanred
    otheraccentcolor = gray,%
    % layercolor must be a colorname from (my-)xcolor; default is black
    layercolor      = black,%
    % mainlanguage must be a babel keyword; default is english
    mainlanguage    = ngerman,%
    % whether a verbose version of the current date should be displayed; default is false
    verbosedate     = true,%
    % whether a verbose version of the current place should be displayed; default is false
    verboseplace    = true,%
    % shift size at the beginning of the text body; default is 225 pt
    refvpos         = 225 pt,%
    % place which is reserved for the addressee field; default is 92 pt
    toaddrheight    = 92pt,%
]{my-letter}

% xcolor is internally loaded.
\definecolor{nf-raisin-black}{rgb}{0.15, 0.15, 0.16}
\definecolor{guardsmanred}{rgb}{0.80, 0.00, 0.00}

\usepackage{mlttr-candidature}

% kantlipsum – Generate sentences in Kant's style.
\RequirePackage{kantlipsum}

\setkomavar{fromlastname}{Kant}
\setkomavar{fromgivenname}{Immanuel}
\setkomavar{fromstreet}{Königsberg, Prussia}
\setkomavar{fromzip}{12345}
\setkomavar{fromemail}{immanuel.kant@example.com}

\setkomavar{place}{Königsberg}
\setkomavar{dateyear}{1804}
\setkomavar{datemonth}{02}
\setkomavar{dateday}{13}
\setkomavar{subject}[]{Application for Professorship in Transcendental Philosophy}

\setkomavar{signature}{Immanuel Kant}

\begin{document}


% Motivational letter
% ENIGMA: Additional environment really necessary?
\begin{letter}{To the Search Committee, Noumenal University}

    \opening{Esteemed Committee,}

    \kant[1-2]

    \closing{With the highest respect,}

    \ps{P.\,S.: <postscriptum>}

    \encl{<attachment1>, ...}

\end{letter}

% CV
\begin{cv}[Curriculum Vitæ]

    % Theses

    \begin{cvblock}{Ph.\,D. Thesis}{Metadata}
        \cvitem{Title}{}{A New Elucidation of the First Principles of Metaphysical Cognition}
        % \cvitem{Advisor}{}{Martin Knutzen}
        \cvitem{Year}{}{1755}
        \cvitem{Description}{}{In this thesis, I endeavored to clarify the first principles of metaphysical cognition, distinguishing the conditions of knowledge from the content of experience.}
    \end{cvblock}

    \begin{cvblock}{Magister Thesis}{Metadata}
        \cvitem{Title}{}{Brief Outline of Certain Meditations on Fire}
        \cvitem{Advisor}{}{Johann Gottfried Teske}
        \cvitem{Year}{}{1755}
        \cvitem{Description}{}{In this thesis, I investigated the nature of fire and its role in natural philosophy, continuing my studies in physics and metaphysics.}
    \end{cvblock}

    % Languages

    \begin{cvblock}{Skills}{}
        \cvitem{Languages}{}{German, Latin, Ancient Greek}
    \end{cvblock}

    \closing{
        The starry heavens above me and the moral law within me.
    }

\end{cv}

\end{document}
